\documentclass[11pt]{article}
\usepackage{amsmath,amssymb,mathtools}
\usepackage{fontspec,unicode-math}
\setmainfont{Latin Modern Roman}
\setmathfont{Latin Modern Math}
\usepackage{xurl,hyperref}
\newcommand{\sref}[2]{\hyperlink{source:#1:#2}{[S#2]}}
\newcommand{\cataloguescope}[1]{\paragraph{Mathematical statement.}#1}
\begin{document}
\section{Take any positive integer, and apply the following process: (1) divide it by two if it's even, multiply by three and add}\label{statement}

\subsection{Definitions and mathematical statement}

\paragraph{Shared notation.}$\mathbb N=\{1,2,\ldots\}$, and $\mathbb Z,\mathbb Q,\mathbb R,\mathbb C$ denote the integers, rationals, reals and complex numbers. Any other conventions are specified locally.


Define $\operatorname{Col}:\mathbb N\to\mathbb N$ by
\[\operatorname{Col}(n)=\begin{cases}n/2,&n\text{ even},\\3n+1,&n\text{ odd}.\end{cases}\]
For $j\ge0$, $\operatorname{Col}^j$ is the $j$-fold iterate, with $\operatorname{Col}^0(n)=n$.

\paragraph{Statement recorded in the supplied source.}
Does every $n\in\mathbb N$ have an integer $j\ge0$ for which $\operatorname{Col}^j(n)=1$? \sref{statement}{2}

\subsection{Short English statement}
Does the repeated halve-or-triple-and-add-one process reach one from every positive starting integer?
\subsection{Sources}
\begin{description}
\item[\hypertarget{source:statement:1}{S1}] ULAM.AI and the UnsolvedMath Contributors, \emph{UnsolvedMath: A Curated Collection of Open Mathematics Problems}, catalogue entry. CC BY 4.0. \url{https://huggingface.co/datasets/ulamai/UnsolvedMath}.
\item[\hypertarget{source:statement:2}{S2}] Terence Tao, \emph{Almost all orbits of the Collatz map attain almost bounded values}, Forum of Mathematics, Pi 10 (2022), e12, Section1.1, definition of the Collatz map and Conjecture1.1, page1. \url{https://arxiv.org/abs/1909.03562}.
\end{description}
\label{end:statement}
\end{document}
